\documentclass[11pt, a4paper]{article}

% Gemeinsame Style-Definitionen (TEXINPUTS muss templates/ enthalten — siehe scripts/build_tex.py)
% bfw-style.tex — gemeinsame Style-Definitionen für alle Handouts/Aufgabenhefte
% Voraussetzung: Kompilation mit xelatex (wegen fontspec)

\usepackage[ngerman]{babel}
\usepackage{geometry}
\geometry{a4paper, top=2.2cm, bottom=2.2cm, left=2.3cm, right=2.3cm}

\usepackage{fontspec}
% Fallback-Kette: Source Sans Pro -> Calibri -> Segoe UI -> Latin Modern Sans
% (Latin Modern ist immer mit MiKTeX vorhanden)
\IfFontExistsTF{Source Sans Pro}{%
  \setmainfont{Source Sans Pro}[Ligatures=TeX]%
}{\IfFontExistsTF{Calibri}{%
  \setmainfont{Calibri}[Ligatures=TeX]%
}{\IfFontExistsTF{Segoe UI}{%
  \setmainfont{Segoe UI}[Ligatures=TeX]%
}{%
  \setmainfont{Latin Modern Sans}[Ligatures=TeX]%
}}}
\IfFontExistsTF{Source Code Pro}{%
  \setmonofont{Source Code Pro}[Scale=0.92]%
}{\IfFontExistsTF{Consolas}{%
  \setmonofont{Consolas}[Scale=0.92]%
}{%
  \setmonofont{Latin Modern Mono}[Scale=0.92]%
}}

\usepackage{xcolor}
% Farbpalette: hell, freundlich, modern
\definecolor{bfwPrimary}{HTML}{0EA5A4}    % Petrol/Türkis - Primär
\definecolor{bfwPrimaryDark}{HTML}{0F766E} % dunkler Petrol für Headings
\definecolor{bfwAccent}{HTML}{F59E0B}     % Amber - Akzent
\definecolor{bfwSuccess}{HTML}{10B981}    % Grün - Erfolg / Lösung
\definecolor{bfwWarn}{HTML}{F97316}       % Orange - Achtung
\definecolor{bfwInk}{HTML}{0F172A}        % fast Schwarz
\definecolor{bfwInkSoft}{HTML}{475569}    % gedämpftes Grau
\definecolor{bfwBgSoft}{HTML}{F8FAFC}     % off-white
\definecolor{bfwBgTeal}{HTML}{ECFEFF}     % helles Türkis
\definecolor{bfwBgAmber}{HTML}{FEF3C7}    % helles Gelb
\definecolor{bfwBgRose}{HTML}{FFE4E6}     % helles Rosé für Eselsbrücken

\usepackage[colorlinks=true, linkcolor=bfwPrimaryDark, urlcolor=bfwPrimaryDark]{hyperref}
\usepackage{enumitem}
\setlist[itemize]{leftmargin=*, itemsep=2pt, topsep=2pt}
\setlist[enumerate]{leftmargin=*, itemsep=2pt, topsep=2pt}

\usepackage[most]{tcolorbox}
\usepackage{booktabs}
\usepackage{tikz}
\usetikzlibrary{decorations.pathmorphing, positioning, shapes.geometric, arrows.meta}

% Merkkästchen — wichtige Definition / Kernpunkt
\newtcolorbox{merksatz}[1][]{%
  enhanced, breakable,
  colback=bfwBgTeal,
  colframe=bfwPrimary,
  boxrule=0.6pt,
  arc=4pt,
  left=10pt, right=10pt, top=8pt, bottom=8pt,
  fonttitle=\bfseries\color{bfwPrimaryDark},
  title={\faLightbulb\ Merksatz},
  #1
}

% Eselsbrücke — Mnemoniks
\newtcolorbox{eselsbruecke}[1][]{%
  enhanced, breakable,
  colback=bfwBgRose,
  colframe=bfwAccent,
  boxrule=0.6pt,
  arc=4pt,
  left=10pt, right=10pt, top=8pt, bottom=8pt,
  fonttitle=\bfseries\color{bfwWarn},
  title={Eselsbrücke},
  #1
}

% Beispiel-Box
\newtcolorbox{beispiel}[1][]{%
  enhanced, breakable,
  colback=bfwBgSoft,
  colframe=bfwInkSoft,
  boxrule=0.4pt,
  arc=3pt,
  left=10pt, right=10pt, top=8pt, bottom=8pt,
  fonttitle=\bfseries\color{bfwInk},
  title={Beispiel},
  #1
}

% Lösungs-Box (für Aufgabenhefte)
\newtcolorbox{loesung}[1][]{%
  enhanced, breakable,
  colback=bfwBgAmber!40,
  colframe=bfwSuccess,
  boxrule=0.6pt,
  arc=3pt,
  left=10pt, right=10pt, top=8pt, bottom=8pt,
  fonttitle=\bfseries\color{bfwSuccess!70!black},
  title={Lösung},
  #1
}

% Glossar-Eintrag
\newcommand{\glossareintrag}[2]{%
  \par\noindent\textbf{\color{bfwPrimaryDark}#1} \enspace #2\par\smallskip
}

% Section-Stil — etwas farbiger als Standard
\usepackage{titlesec}
\titleformat{\section}
  {\Large\bfseries\color{bfwPrimaryDark}}
  {\thesection}{0.6em}{}
\titleformat{\subsection}
  {\large\bfseries\color{bfwPrimary}}
  {\thesubsection}{0.5em}{}
\titleformat{\subsubsection}
  {\normalsize\bfseries\color{bfwInk}}
  {\thesubsubsection}{0.4em}{}

% TikZ-Stil "skizzenhaft" — etwas weicher als Rough.js, aber stabil
\tikzset{
  roughbox/.style = {
    draw=bfwPrimaryDark, thick, fill=bfwBgTeal,
    rounded corners=4pt, inner sep=6pt
  },
  roughaccent/.style = {
    draw=bfwAccent, thick, fill=bfwBgAmber,
    rounded corners=4pt, inner sep=6pt
  },
  roughline/.style = {
    draw=bfwInkSoft, thick, -{Stealth[length=2.5mm]}
  }
}

% Bullet-Symbole (bewusst kein FontAwesome — vermeidet Font-Installations-Probleme)
\newcommand{\faLightbulb}{$\bullet$}
\newcommand{\faBookmark}{$\bullet$}

% Header/Footer
\usepackage{fancyhdr}
\pagestyle{fancy}
\fancyhf{}
\renewcommand{\headrulewidth}{0pt}
\fancyfoot[L]{\small\color{bfwInkSoft}\thepage}
\fancyfoot[R]{\small\color{bfwInkSoft} BFW M-064 Beschaffung im Büromanagement}


\title{\vspace*{-1.5cm}\Huge\bfseries\color{bfwPrimaryDark}Tag~13: Ablage \& Ablagesysteme\\[0.4em]
  \large\color{bfwInkSoft}\textnormal{Schriftgut ordnen, aufbewahren und wiederfinden}}
\author{\textsf{Büroprozesse gestalten und Arbeitsvorgänge organisieren \textperiodcentered{} M-067}\\
\textsf{\small\copyright\ Can Siebert\,\textperiodcentered\,2026}}
\date{15.07.2026}

\begin{document}
\maketitle
\thispagestyle{empty}

\vspace{1em}
\begin{merksatz}[title={\bfseries Worum geht es heute?}]
Eine Akte ist nur so viel wert, wie schnell man sie wiederfindet. Wer Schriftgut
planlos ablegt, verliert Zeit, Unterlagen und am Ende Geld. In diesem Handout lernen Sie,
\emph{wo} Unterlagen aufbewahrt werden (Ablageorte), \emph{wie} sie aufbewahrt werden
(Ablagetechnik und Registraturformen) und \emph{nach welcher Regel} sie geordnet werden
(Ordnungssysteme). Dazu kommen Ordnungshilfsmittel wie der Aktenplan und ein Blick auf
die Kosten der Ablage. Das ist die Grundlage für die digitale Dokumentenverwaltung am
Folgetag und ein wiederkehrendes Thema der IHK-Abschlussprüfung.
\end{merksatz}

\tableofcontents
\vspace{1em}
\hrule
\vspace{1em}

\section{Ablageorte --- wo wird abgelegt?}

Bei der Entscheidung, an welchem Ort die Unterlagen aufbewahrt werden, sollte das
Unternehmen einen einfachen Grundsatz beachten: \emph{Je häufiger ein Mitarbeiter auf
die Schriftstücke zurückgreifen muss, umso näher an seinem Arbeitsplatz sollten sie
aufbewahrt werden.} Daraus ergeben sich drei mögliche Ablageorte.

\begin{enumerate}
  \item \textbf{Arbeitsplatzablage} --- Schriftstücke, die der Mitarbeiter aktuell
    bearbeitet oder auf die er häufig und schnell zurückgreifen muss, bewahrt er direkt
    an seinem Arbeitsplatz auf. Er sorgt dabei selbst für eine geeignete Form der Ablage.
  \item \textbf{Abteilungsablage} --- Werden Schriftstücke von mehreren Mitarbeitern
    gleichzeitig bearbeitet oder verwaltet, sollten sie für alle in der Abteilung gut
    erreichbar abgelegt werden. Die Ablage sollte für alle nachvollziehbar und
    verständlich organisiert sein --- hilfreich ist hier ein Aktenplan.
  \item \textbf{Zentralregistratur / Archiv} --- Hier werden die Unterlagen aus den
    verschiedenen Bereichen des Unternehmens aufbewahrt. Die Zentralregistratur entlastet
    die Abteilungen von Schriftstücken, die nicht mehr benötigt werden, aber noch
    aufbewahrt werden müssen (z.\,B.\ wegen der gesetzlichen Aufbewahrungspflicht). Auch
    Unterlagen mit \emph{Dauerwert} werden hier archiviert.
\end{enumerate}

\begin{merksatz}
Zugriffshäufigkeit entscheidet über den Ort: Was täglich gebraucht wird, gehört an den
Arbeitsplatz; was selten oder nur aus rechtlichen Gründen benötigt wird, wandert ins
Archiv.
\end{merksatz}

\section{Ablagetechnik und Aktenarten}

\subsection{Loseblattablage und geheftete Ablage}

Bei der \textbf{Loseblattablage} werden die Schriftstücke ungeheftet, also lose, in die
Schriftgutbehälter (z.\,B.\ Mappen, Sichthüllen) eingelegt. Vorteilhaft ist der
\emph{Zeitgewinn} bei der Ablage, da nicht gelocht werden muss. Nachteilig ist die Gefahr,
dass Unterlagen verloren gehen; bei umfangreichen Akten können sich längere Suchzeiten
ergeben. Geeignet ist diese Technik für Schriftstücke, die im Arbeitsprozess
weitergegeben werden, nicht zu umfangreich sind und schnell zur Hand sein müssen.

Bei der \textbf{gehefteten Ablage} werden die Schriftstücke gelocht und in die
Schriftgutbehälter (z.\,B.\ Ordner, Hefter) eingeheftet. Diese Art erfordert mehr Zeit
(Lochen vor dem Abheften), bietet aber eine \emph{sichere Aufbewahrung} und ein
\emph{schnelles Wiederfinden}. Geeignet ist sie für eine größere Anzahl gleichartiger
Schriftstücke, die nach einem einheitlichen Kriterium fortlaufend abgelegt werden ---
z.\,B.\ Aufträge nach Auftragsnummer.

\subsection{Einzelakte und Sammelakte}

\begin{itemize}
  \item \textbf{Einzelakte} --- In ihr werden die Schriftstücke zu \emph{einem} Vorgang
    oder Fall abgelegt. Beispiel: Die Personalakte eines Mitarbeiters.
  \item \textbf{Sammelakte} --- In ihr werden Schriftstücke zu verschiedenen, aber
    \emph{gleichartigen} Vorgängen abgelegt. Beispiel: Lieferscheine oder Rechnungen.
\end{itemize}

\begin{eselsbruecke}
\textbf{Einzel = ein Fall, Sammel = viele gleiche.} Die Personalakte (ein Mensch) ist
eine Einzelakte; der Ordner „Eingangsrechnungen” (viele gleichartige Belege) ist eine
Sammelakte.
\end{eselsbruecke}

\section{Arten der Registratur}

Es lassen sich drei grundlegende \emph{Aufbewahrungsformen} unterscheiden:
\textbf{liegende}, \textbf{stehende} und \textbf{hängende} Registratur. Jeder
Aufbewahrungsform lassen sich unterschiedliche Registraturformen und Schriftgutbehälter
zuordnen.

\subsection{Liegende Registratur}

Bei der \textbf{Flachablage} werden die Schriftstücke flach übereinandergelegt.
Schriftgutbehälter sind Aktendeckel, Schnellhefter und Jurismappe. Sie eignet sich für
Akten, die aktuell im laufenden Arbeitsprozess oder nur selten benötigt werden.
\emph{Vorteil}: geringe Kosten für die Schriftgutbehälter. \emph{Nachteil}: schlechte
Übersicht, die Bearbeitung ist sehr umständlich.

\subsection{Stehende Registratur}

\textbf{Ordnerregistratur} --- Schriftgutbehälter ist der Ordner. Geeignet für
umfangreiche Sammelakten (z.\,B.\ Rechnungskopien, Lieferscheine). \emph{Vorteile}: sehr
übersichtlich durch die beschriftbaren Ordnerrücken, schneller Zugriff, geringe Kosten,
Schriftstücke werden sicher abgeheftet. \emph{Nachteile}: zeitaufwendig (Lochen und
Abheften); bei unregelmäßigem Aktenanfall muss Platz frei gelassen werden, was die
Ablageart unflexibel macht.

\textbf{Stehsammlerregistratur} --- Schriftgutbehälter sind Einstellmappen, Stehsammler
sowie Kassetten und Archivschachteln. Geeignet für Sammelakten und nicht allzu
umfangreiche Einzelakten. \emph{Vorteile}: schneller Zugriff, gute Raumausnutzung durch
die Schriftgutbehälter. \emph{Nachteile}: durch die feststehende Bodenbreite der
Kassetten fehlt die Flexibilität; bei unregelmäßigem Aktenanfall muss Platz frei
gehalten werden.

\subsection{Hängende Registratur}

\textbf{Hängeregistratur} --- Schriftgutbehälter sind Hängemappen, Hängetaschen,
Hängeordner, Hängehefter und Hängesammler. Geeignet für Handakten am Arbeitsplatz und
größere Akten der Abteilungsablage (auch ungeheftet ablegbar). \emph{Vorteile}: sehr
übersichtlich, schneller Zugriff (bei Verwendung von Mappen und Taschen), sehr flexibel.
\emph{Nachteile}: höhere Materialkosten durch Schriftgutbehälter und Registraturmöbel,
größerer Raumbedarf, da nur bis zur Sichthöhe abgelegt werden kann.

\textbf{Pendelregistratur} --- Schriftgutbehälter sind Pendelmappen, Pendeltaschen,
Pendelhefter und Pendelsammler. Geeignet für eine große Anzahl von Einzelakten in der
Abteilungsablage. \emph{Vorteile}: bessere Ausnutzung der Stellfläche (keine Auszüge
nötig, Raum bis Griffhöhe nutzbar), Schriftgutbehälter leicht verschiebbar und daher
sehr flexibel, relativ günstige Anschaffungskosten der Registraturmöbel.
\emph{Nachteile}: längere Zugriffszeit, da der gesamte Schriftgutbehälter entnommen
werden muss; weniger übersichtlich, da man nur die Schmalseite sieht.

\section{Ordnungssysteme}

In einer gut organisierten Registratur sollen alle Mitarbeiter abgelegte Schriftstücke
schnell und ohne Rückfragen finden. Das setzt ein \emph{allgemein anerkanntes und allen
bekanntes Ordnungssystem} voraus. Folgende Ordnungssysteme lassen sich unterscheiden:

\begin{description}[leftmargin=1.2cm, style=nextline]
  \item[Alphabetische Ordnung] Ablage in der Reihenfolge der Ordnungswörter (z.\,B.\
    Kunden- oder Lieferantenname). Maßgeblich ist nach \emph{DIN 5007} die Buchstabenfolge
    des Alphabets.
  \item[Numerische Ordnung] Die Schriftstücke werden mit einer Nummer versehen und
    entsprechend abgelegt (z.\,B.\ Lieferschein-Nr., Kunden-Nr., Lieferanten-Nr.).
  \item[Alphanumerische Ordnung] Kombination aus alphabetischer und numerischer Ordnung;
    findet häufig bei der mnemotechnischen Ordnung Anwendung.
  \item[Mnemotechnische Ordnung] Vom griechischen \emph{Mneme} = Gedächtnis. Es werden
    Gedächtnishilfen genutzt: Für Ausgangsrechnungen z.\,B.\ das Kürzel \emph{AR}, für
    Eingangsrechnungen \emph{ER}; die Rechnungsnummer dient der genauen Einordnung.
  \item[Chronologische Ordnung] Ablage nach der zeitlichen Reihenfolge (Tage, Wochen,
    Monate oder Jahre). Das aktuellste Schriftstück liegt ganz vorne oder ganz hinten.
    Bilanzen und Steuerbescheide werden z.\,B.\ nach Jahren abgelegt.
  \item[Geografische Ordnung] (ergänzende Praxis) Ablage nach räumlichen Kriterien wie
    Land, Region, Postleitzahl oder Vertriebsgebiet --- nützlich z.\,B.\ im Außendienst.
\end{description}

\begin{beispiel}
Alphabetische Reihenfolge nach DIN 5007: \emph{Kaiser, Krähmer, Kuhn, Kurz, Larsson,
Ludwig.} Bei gleichem Anfangsbuchstaben entscheidet der zweite, dritte usw.\ Buchstabe;
Umlaute werden in der Wörterbuch-Variante wie der Grundbuchstabe behandelt (ä = a),
„ß” wird wie „ss” eingeordnet.
\end{beispiel}

\subsection{Vergleich der Ordnungssysteme}

\begin{center}
\begin{tabular}{p{2.9cm} p{5.2cm} p{4.4cm}}
\toprule
\textbf{Ordnungssystem} & \textbf{Ordnungskriterium} & \textbf{Typisches Beispiel}\\
\midrule
Alphabetisch & Buchstabenfolge der Ordnungswörter (DIN 5007) & Kundenkartei nach Namen\\
Numerisch & fortlaufende Nummer & Belege nach Lieferschein-Nr.\\
Alphanumerisch & Kombination Buchstabe + Zahl & Kunden-Nr.\ mit Namenskürzel\\
Mnemotechnisch & Gedächtnishilfe (Kürzel) + Zahl & AR / ER + Rechnungsnummer\\
Chronologisch & zeitliche Reihenfolge & Bilanzen nach Jahren\\
Geografisch & räumliches Kriterium & Akten nach Vertriebsgebiet\\
\bottomrule
\end{tabular}
\end{center}

\begin{eselsbruecke}
\textbf{„A-N-A-M-C-G”} --- \emph{A}lphabetisch, \emph{N}umerisch, \emph{A}lphanumerisch,
\emph{M}nemotechnisch, \emph{C}hronologisch, \emph{G}eografisch. So behalten Sie die
sechs Ordnungssysteme im Kopf.
\end{eselsbruecke}

\section{Ordnungshilfsmittel}

Um die Arbeit mit abgelegtem Schriftgut zu erleichtern, steht eine Vielzahl von
\emph{Ordnungshilfsmitteln} zur Verfügung:

\begin{itemize}
  \item \textbf{Register- und Trennblätter} helfen, Übersicht in Ordnern herzustellen;
    verfügbar mit alphabetischen und numerischen Ordnungsreihen.
  \item \textbf{Leitkarten} ermöglichen bei Hängemappen oder Karteikästen eine
    Unterteilung der Schriftstücke in Gruppen.
  \item \textbf{Reiter} bestehen aus Metall oder Kunststoff (Fenster- oder
    Vollsichtreiter), tragen Inhaltsschilder und werden auf die Schriftgutbehälter
    gesteckt.
  \item \textbf{Wiedervorlagemappe} --- siehe unten.
  \item \textbf{Aktenplan} --- siehe unten.
\end{itemize}

\subsection{Wiedervorlagemappe}

In einer \textbf{Wiedervorlagemappe} werden alle Aufgaben und Vorgänge abgelegt, die
nicht sofort bearbeitet oder erledigt werden --- also noch zu bearbeitende Schriftstücke
oder schwebende Vorgänge. So werden gestapelte Schriftstücke auf dem Schreibtisch
vermieden. Sie ist ein wichtiges organisatorisches Hilfsmittel bei der Terminplanung:
Vorgänge werden termingerecht bearbeitet, und die Einhaltung der Termine kann überprüft
werden. Die Mappe ist entweder in zwölf Fächer für die einzelnen Monate oder in 31 Fächer
für die Monatstage unterteilt. Die Trennblätter weisen meist Sichtlöcher auf, damit man
erkennt, ob sich Schriftstücke in der Mappe befinden.

\subsection{Aktenplan und Dezimalklassifikation}

Der \textbf{Aktenplan} (Informationsstrukturplan) ist der einheitliche Rahmen in einem
Unternehmen, nach dem Schriftstücke abgelegt werden. Er enthält die systematische und
logische Ordnung für die Schriftgutablage. Mithilfe von Sachgebieten (= Hauptgruppen),
Untergruppen und Stichworten wird die Struktur bestimmt, anhand derer die Unterlagen in
das betriebliche Archiv eingeordnet werden. Da einheitliche Regeln gelten, fällt die
Einordnung leichter und die Dokumente werden von allen Mitarbeitern schnell und sicher
gefunden. Häufig ist der Aktenplan als \emph{Dezimalklassifikation} aufgebaut: Die
Hauptgruppen werden mit den Ziffern 0--9 nummeriert, Untergruppen und Stichworte durch
weitere Ziffern verfeinert (z.\,B.\ 5 Personal $\rightarrow$ 51 Personalverwaltung
$\rightarrow$ 514 Urlaubsanträge).

\begin{beispiel}
Auszug aus dem Aktenplan der \emph{Duisdorfer BüroKonzept KG}: 0~Geschäftsführung,
1~Anlagen, 2~Beschaffung, 3~Vertrieb, 4~Kfm.\ Verwaltung, 5~Personal (51~Personal\-verwaltung,
52~Personalvergütung, 53~Vertragswesen), 6~Technischer Bereich, 7~Öffentlichkeitsarbeit,
8~Qualitätssicherung, 9~frei. So liegt z.\,B.\ ein Urlaubsantrag unter 514.
\end{beispiel}

\subsection{Elektronische Dateiablage}

Bei einer \textbf{elektronischen Dateiablage} sollten Texte und Daten --- falls kein
Datenmanagementsystem genutzt wird --- nach einem einheitlichen Dateiablagesystem
gespeichert werden. Dabei sollten die Ordnerstrukturen der elektronischen Ablage und der
Papierablage aufeinander abgestimmt sein. Auch die Dateibenennung sollte einheitlich
geregelt sein (vgl.\ die DIN-Schreib- und Gestaltungsregeln für die Text- und
Informationsverarbeitung). Das ist die Brücke zur digitalen Dokumentenverwaltung (DMS)
am Folgetag.

\section{Kosten der Ablage}

Die Ablage verursacht Kosten, die bei der Wahl des Systems mitbedacht werden müssen:

\begin{itemize}
  \item \textbf{Materialkosten} --- Schriftgutbehälter und Registraturmöbel (gering bei
    der Flachablage und der Pendelregistratur, höher bei der Hängeregistratur).
  \item \textbf{Raumkosten} --- benötigte Stell- bzw.\ Stapelfläche (die
    Hängeregistratur braucht mehr Raum, da nur bis Sichthöhe abgelegt wird; die
    Pendelregistratur nutzt den Raum bis zur Griffhöhe besser aus).
  \item \textbf{Arbeitszeit} --- Lochen, Abheften und Suchen (die geheftete Ablage und
    die Ordnerregistratur sind zeitaufwendiger als die Loseblattablage).
\end{itemize}

\begin{merksatz}
Die günstigste Registraturform ist nicht automatisch die wirtschaftlichste: Geringe
Materialkosten nützen wenig, wenn das Suchen und Bearbeiten viel Arbeitszeit kostet.
Material-, Raum- und Zeitkosten sind immer zusammen zu betrachten.
\end{merksatz}

\section{Zusammenfassung}

\begin{enumerate}
  \item \textbf{Ablageorte}: Arbeitsplatz-, Abteilungsablage und Zentralregistratur/Archiv --- entscheidend ist die Zugriffshäufigkeit.
  \item \textbf{Ablagetechnik}: Loseblattablage (schnell, aber unsicher) und geheftete Ablage (sicher, aber zeitaufwendig); Akten als Einzel- oder Sammelakte.
  \item \textbf{Registraturformen}: liegend (Flachablage), stehend (Ordner-/Stehsammlerregistratur), hängend (Hänge-/Pendelregistratur) --- jeweils mit eigenen Vor- und Nachteilen.
  \item \textbf{Ordnungssysteme}: alphabetisch (DIN 5007), numerisch, alphanumerisch, mnemotechnisch, chronologisch und geografisch.
  \item \textbf{Ordnungshilfsmittel}: Register, Reiter, Leitkarten, Wiedervorlagemappe und der Aktenplan (oft als Dezimalklassifikation 0--9).
  \item \textbf{Kosten der Ablage}: Material-, Raum- und Arbeitszeitkosten gemeinsam abwägen --- die Brücke zur digitalen Ablage (DMS) folgt an Tag~14.
\end{enumerate}

\newpage

\section*{Glossar}
\addcontentsline{toc}{section}{Glossar}

\glossareintrag{Ablageort}{Ort, an dem Unterlagen aufbewahrt werden; nach Zugriffshäufigkeit Arbeitsplatz-, Abteilungsablage oder Zentralregistratur/Archiv.}

\glossareintrag{Zentralregistratur}{Zentrale Aufbewahrung von Unterlagen aus allen Bereichen; entlastet Abteilungen und archiviert Unterlagen mit Dauerwert.}

\glossareintrag{Loseblattablage}{Ungeheftete Ablage in Mappen oder Sichthüllen; zeitsparend, aber mit Verlustgefahr.}

\glossareintrag{Geheftete Ablage}{Gelochte Ablage in Ordnern oder Heftern; sichere Aufbewahrung und schnelles Wiederfinden.}

\glossareintrag{Einzelakte}{Akte mit den Schriftstücken zu einem Vorgang oder Fall, z.\,B.\ Personalakte.}

\glossareintrag{Sammelakte}{Akte mit Schriftstücken zu verschiedenen, aber gleichartigen Vorgängen, z.\,B.\ Rechnungen.}

\glossareintrag{Flachablage}{Liegende Registratur; Schriftstücke werden flach übereinandergelegt (Aktendeckel, Schnellhefter, Jurismappe).}

\glossareintrag{Ordnerregistratur}{Stehende Registratur mit Ordnern; übersichtlich durch beschriftbare Rücken, aber zeitaufwendig.}

\glossareintrag{Stehsammlerregistratur}{Stehende Registratur mit Einstellmappen, Stehsammlern, Kassetten und Archivschachteln.}

\glossareintrag{Hängeregistratur}{Hängende Registratur mit Hängemappen und -taschen; übersichtlich und flexibel, aber raum- und materialintensiv.}

\glossareintrag{Pendelregistratur}{Hängende Registratur mit Pendelmappen; gute Raumnutzung bis Griffhöhe, aber längere Zugriffszeit.}

\glossareintrag{Schriftgutbehälter}{Behältnis zur Aufbewahrung von Schriftstücken, z.\,B.\ Ordner, Mappe, Hefter, Stehsammler.}

\glossareintrag{DIN 5007}{Norm, die die alphabetische Sortierung (Buchstabenfolge des Alphabets, Behandlung von Umlauten und „ß”) regelt.}

\glossareintrag{Mnemotechnische Ordnung}{Ordnung mithilfe von Gedächtnishilfen/Kürzeln (z.\,B.\ AR für Ausgangsrechnung), ergänzt durch Nummern.}

\glossareintrag{Chronologische Ordnung}{Ablage nach zeitlicher Reihenfolge (Tage, Wochen, Monate, Jahre).}

\glossareintrag{Wiedervorlagemappe}{Mappe (12 oder 31 Fächer) für noch zu bearbeitende Vorgänge; Hilfsmittel der Terminplanung.}

\glossareintrag{Aktenplan}{Einheitlicher Informationsstrukturplan mit Hauptgruppen, Untergruppen und Stichworten für die Schriftgutablage.}

\glossareintrag{Dezimalklassifikation}{Aufbau des Aktenplans mit den Ziffern 0--9 für Haupt- und Untergruppen.}

\end{document}
